% explainer-source-hash: f4056f58dd7b4a388ee50f223c77b4315cabd8562bb0f0a0f393b8a5e9d6e423
% explainer-source-commit: 2b8bb12994f7dc6f75fe0d99b2cff9e1de273266
% explainer-generated: 2026-07-16
% Regenerate with /explain-all (see WIRING.md). Do not edit this stamp by hand:
% it is how the toolchain knows whether this document still matches the code.
\documentclass[11pt,a4paper]{article}
\input{preamble}

\begin{document}

\begin{titlepage}
\centering
\vspace*{2.5cm}
{\Huge\bfseries\color{inkblue} Agency Blog SaaS: Demo\par}
\vspace{0.6cm}
{\Large\color{accent} Deep Dive\par}
\vspace{1.2cm}
{\large A scripted, fabricated product tour of a private production system\par}
\vspace{2.5cm}
\begin{minipage}{0.86\textwidth}
\begin{honest}{Read this before anything else}
This document describes \file{agency-blog-saas-demo}: a small static web page that
\textbf{imitates} a product. It is not the product.

The real system, \file{agency-blog-saas}, is a production Django application serving
paying clients of a real agency. Its source is private and stays private, because it is
the agency's live product rather than a codebase that could be scrubbed and published.
This demo exists so that the product's screens can be clicked through by someone who will
never be given the real repository.

Almost everything the demo displays is \textbf{fabricated}: the credit balance, the client,
the client's updates, the generation process itself. Nothing is computed. Nothing is sent
anywhere. The one exception is the blog library, which reproduces real, already-public
articles from a real client's public website, and this document is careful to say which is
which every time it matters.
\end{honest}
\end{minipage}
\vfill
{\small Built from the project source as tracked in its git repository.\par}
\end{titlepage}

\tableofcontents
\newpage

\section{What this project is}

\subsection{The one-paragraph version}

\file{agency-blog-saas-demo} is a single web page, written in plain HTML, CSS and
JavaScript, with no build step, no framework, no server of its own and no network calls at
runtime. It presents four tabs that mirror the main screens of a private Django product:
an AI Blog Generator, a Blog Library, Subscription Packages, and Client Updates, plus a
fifth hidden panel that reads a full article in place. Clicking \code{Generate post} runs a
pre-written timeline of status messages, subtracts a number from a variable, and copies a
post that was already stored in the page into the library with a countdown on it. That is
the whole program.

\begin{keyidea}{The distinction that governs this entire document}
There are two separate things with similar names.

\textbf{The real system} (\file{agency-blog-saas}) is a production Django SaaS behind
digitalise.agency. It has a database, user accounts, real billing, real AI generation, PDF
export, and real clients. \textbf{You cannot see it}, and neither can a hiring panel.

\textbf{This demo} (\file{agency-blog-saas-demo}) is roughly 1{,}200 lines of static files,
of which the large majority is copied article text. It has no database, no accounts, no
billing, no AI, and no clients. \textbf{It only knows how to draw pictures of those things.}

Every claim in this document is about the demo unless it explicitly says otherwise. Where
the demo makes a claim about the real system, this document treats that claim as a claim,
not as a verified fact, because the code here cannot verify it.
\end{keyidea}

\subsection{Why a demo repository exists at all}

A portfolio has a structural problem: the most substantial work is often the work that
cannot be shown. An agency's live product is commercially sensitive in a way that no amount
of secret-scrubbing fixes, because the sensitive thing is not a password in a config file,
it is the product. Deleting the credentials from a live product's source and publishing the
rest still publishes the product.

There are three usual responses, and it is worth naming why this project chose the third.

\begin{enumerate}
\item \textbf{Publish nothing.} Honest, but the reader has only a written claim to go on.
\item \textbf{Publish screenshots.} Better, but a screenshot cannot be clicked, and a
      reader cannot tell a screenshot of a working product from a screenshot of a mockup.
\item \textbf{Publish a fabricated, clearly-labelled reimplementation of the screens.}
      The reader gets something to click. The cost is that what they are clicking is a
      puppet, and the whole design problem becomes making sure nobody mistakes the puppet
      for the real thing.
\end{enumerate}

This project is option three, and most of the interesting decisions in it are about
managing the cost of that choice rather than about the code.

\begin{jargon}{Static site}
A website made of files that a server hands over exactly as they are stored: HTML, CSS,
images, JavaScript. Nothing is generated per visitor. The opposite is a \emph{dynamic}
site, where a program on the server (Django, for example) builds a different page for each
request by reading a database. The real system is dynamic. This demo is static, which is
precisely why it can pretend to be dynamic only by scripting the pretence in the browser.
\end{jargon}

\subsection{What is honestly demonstrated, and what is not}

This matters enough to be a table rather than a paragraph.

\begin{table}[H]
\centering
\small
\begin{tabular}{@{}p{0.30\textwidth}p{0.30\textwidth}p{0.32\textwidth}@{}}
\toprule
\textbf{Product capability} & \textbf{What the demo shows} & \textbf{What that proves} \\
\midrule
AI blog generation &
A button, a form, three status messages on timers, and a pre-stored post appearing. &
That a generation screen exists and roughly what it feels like. \emph{Nothing} about
prompting, model choice, queueing, or output quality. \\
\addlinespace
Credits and metering &
A number in the corner that decreases by 100. &
The concept. Not the ledger, not the refund-on-failure logic, not the race conditions
that make metering hard. \\
\addlinespace
Subscriptions &
Three price cards and a redemption box that always says it is a demo. &
The tier structure and price points, which are real and public. Not payment, not
provisioning, not renewal. \\
\addlinespace
Per-client editor roles &
Nothing. There is no login and no notion of a user. &
Nothing. This capability is named in the demo's own prose but is absent from the demo. \\
\addlinespace
PDF export &
Nothing. There is no export button. &
Nothing. Also named in the prose, also absent. \\
\addlinespace
Client updates &
A three-item list for an openly invented company. &
The screen's shape only. \\
\addlinespace
Published output quality &
Fifteen genuine, full-length articles from a real client's live blog. &
This is the one place where something real is on display. See section~\ref{sec:tiers}. \\
\bottomrule
\end{tabular}
\caption{Capability by capability: the demo is a picture of a product, and two of the
capabilities it advertises in its own hero text are not pictured at all.}
\end{table}

\begin{honest}{The demo advertises two features it does not contain}
The page's own opening paragraph tells the visitor that the real system does ``AI blog
generation, credits and subscriptions, per-client editor roles, and PDF export''. The first
two have tabs. \textbf{Per-client editor roles and PDF export have no representation
anywhere in the demo}: no login, no role switcher, no export control, no mention on any
tab.

This is not dishonest, because the sentence is explicitly describing the real system rather
than the demo, and the very next sentence says the page is a scripted tour. But it is a
real gap between the advertised feature list and the tour: a visitor who reads the list and
then clicks all four tabs looking for roles and export will not find them, and will not be
told why. Either the tour should cover them or the hero should not lead with them.
\end{honest}

\section{Running it}

\begin{lstlisting}[caption={There is no setup. These are the only two ways to run it.}]
# 1. Open the file directly.
xdg-open index.html          # or double-click it

# 2. Or serve the folder with any static server.
python3 -m http.server 8000  # then visit http://localhost:8000
\end{lstlisting}

Both work identically, because the page never makes a request that a \code{file://} origin
would block. There is no \code{npm install}, no compiler, no environment file and no
secrets, and this is a deliberate property rather than an accident: see
section~\ref{sec:selfcontained}.

\file{netlify.toml} is the entire deployment configuration:

\begin{lstlisting}[caption={netlify.toml, in full}]
[build]
  publish = "."
\end{lstlisting}

\begin{jargon}{Netlify, and what \code{publish = "."} means}
Netlify is a hosting service that watches a git repository and serves its contents as a
website. Most projects need a \emph{build command} first, turning source into deployable
files. This project has no build, so the configuration says only ``publish the folder as
it is''. Two lines of deployment configuration is a direct consequence of the
no-build-step decision.
\end{jargon}

\section{The layout of the repository}

\begin{figure}[H]
\centering
\begin{forest} dirtree
[agency-blog-saas-demo
  [index.html {\rmfamily\itshape\scriptsize\quad 137 lines: all five panels, nav, footer}]
  [demo-data.js {\rmfamily\itshape\scriptsize\quad 375 lines, but $\approx$150\,KB: the data}]
  [app.js {\rmfamily\itshape\scriptsize\quad 216 lines: every behaviour on the page}]
  [styles.css {\rmfamily\itshape\scriptsize\quad 285 lines: the borrowed design system}]
  [netlify.toml {\rmfamily\itshape\scriptsize\quad 2 lines}]
  [favicon.svg]
  [favicon.png {\rmfamily\itshape\scriptsize\quad fallback for browsers without SVG icons}]
  [assets
    [fonts
      [space-grotesk.woff2]
    ]
    [posts
      [CREDITS.md {\rmfamily\itshape\scriptsize\quad provenance of the real content}]
      [{15 x *.webp} {\rmfamily\itshape\scriptsize\quad one hero image per real post}]
    ]
  ]
  [docs
    [explainers
      [{deep-dive.tex, brief.tex, and their PDFs}]
    ]
  ]
  [README.md]
  [LICENSE {\rmfamily\itshape\scriptsize\quad PolyForm Strict 1.0.0}]
]
\end{forest}
\caption{The whole project. Four files carry the entire program; the 2.5\,MB of weight is
almost all images and one font.}
\end{figure}

The file count is small enough that the architecture is really just an answer to one
question: what goes in which of the four files? The answer is a clean and conventional
split, and it is worth stating because it is one of the things a reader can actually judge
here.

\begin{itemize}
\item \file{index.html} holds \emph{structure only}: five \code{<section>} panels, each
      mostly empty, waiting to be filled.
\item \file{demo-data.js} holds \emph{data only}, and defines three global constants.
\item \file{app.js} holds \emph{behaviour only}, and reads those constants.
\item \file{styles.css} holds \emph{appearance only}.
\end{itemize}

\begin{honest}{The split is real but not absolute}
Three of the panels in \file{index.html} carry inline \code{style="..."} attributes
(\code{color:var(--ink-soft);font-size:14px} and similar) rather than classes, so a little
appearance has leaked into the structure file. It is cosmetic and it is four attributes,
but the separation the project otherwise keeps carefully is not perfectly kept.
\end{honest}

\section{The three tiers of truth}
\label{sec:tiers}

This is the most important idea in the project, and it is not a coding idea.

The data on the page is not uniformly fake. It comes in three tiers with genuinely
different provenance, and the design of the page is largely an attempt to keep a visitor
from confusing them.

\begin{figure}[H]
\centering
\resizebox{\textwidth}{!}{
\begin{tikzpicture}[node distance=6mm]

\node[box, fill=warnred!10, draw=warnred, text width=4.1cm, minimum height=3.4cm,
      anchor=north west] (t1) at (0,0)
  {\textbf{Tier 1: fabricated}\\[3pt]
   \scriptsize Invented for this demo.\\ Nothing behind it.\\[4pt]
   \scriptsize $\bullet$ 1{,}200 credit balance\\
   \scriptsize $\bullet$ ``Aurora Fitness Co.''\\
   \scriptsize $\bullet$ its three updates\\
   \scriptsize $\bullet$ the generation timeline};

\node[box, fill=softgrey, text width=4.1cm, minimum height=3.4cm,
      anchor=north west] (t2) at (4.7,0)
  {\textbf{Tier 2: real and public,}\\\textbf{not client-specific}\\[3pt]
   \scriptsize Real facts, already published.\\[4pt]
   \scriptsize $\bullet$ the three price tiers\\
   \scriptsize $\bullet$ the tech stack line\\
   \scriptsize $\bullet$ the test/coverage line\\
   \scriptsize $\bullet$ davonex.com's design tokens};

\node[box, fill=okgreen!10, draw=okgreen, text width=4.1cm, minimum height=3.4cm,
      anchor=north west] (t3) at (9.4,0)
  {\textbf{Tier 3: real client work,}\\\textbf{already public}\\[3pt]
   \scriptsize Genuine shipped output.\\[4pt]
   \scriptsize $\bullet$ 15 davonex.com articles\\
   \scriptsize $\bullet$ their hero images\\
   \scriptsize $\bullet$ their references\\
   \scriptsize $\bullet$ their real dates};

\node[extbox, text width=4.1cm, anchor=north west] (x1) at (0,-3.8)
  {\scriptsize Labelled by the red banner\\ at the top of every page,\\ and by naming the
   client\\ ``(invented, illustrative)''\\ in the data itself.};
\node[extbox, text width=4.1cm, anchor=north west] (x2) at (4.7,-3.8)
  {\scriptsize Safe because it was\\ already public before this\\ demo existed. Reused\\
   verbatim rather than\\ approximated.};
\node[extbox, text width=4.1cm, anchor=north west] (x3) at (9.4,-3.8)
  {\scriptsize Labelled by a provenance\\ note and a link back to\\ the original URL at the\\
   foot of every article.};

\draw[flow] (t1.south) -- (x1.north);
\draw[flow] (t2.south) -- (x2.north);
\draw[flow] (t3.south) -- (x3.north);

\node[box, draw=accent, fill=white, text width=13.5cm, anchor=north west] (rule) at (0,-6.4)
  {\textbf{The rule that ties them together:} anything fabricated is duplicated from
   something real and stamped \code{Demo}, so the visitor is never asked to guess. A demo
   card in the library is always a copy of a Tier-3 post, never an invention, and it always
   carries the ribbon and the countdown.};
\end{tikzpicture}}
\caption{Three provenances on one page, and the labelling that keeps them apart.}
\end{figure}

\subsection{Tier 3 is the unusual choice}

Most fabricated demos invent their content library too: a page of lorem-ipsum blog posts
with plausible titles. This one does not. It ships all fifteen real, live articles from
\file{davonex.com/blog}, an actual client, with their real titles, real publication dates,
real read times, real hero images, real full article bodies (tables, headings, inline
citations and all) and their real reference lists. Verified by counting them in the source:
15 posts, 15 unique slugs, 15 unique titles, 177 reference entries in total, and roughly
154\,KB of article HTML.

The project's \file{README.md} argues that this is a \emph{stronger} honesty position than
inventing a library, and the argument is sound: a reader can see genuine shipped work and a
working interaction in one place, with a link back to the original on every article so any
claim is checkable in one click.

\begin{honest}{But Tier 3 comes from a different private system than the one being demoed}
\file{assets/posts/CREDITS.md} states that davonex.com is ``an actual configured client of
the real \file{blog-posting-pipeline} production system''. \file{blog-posting-pipeline} is a
\emph{different} private project from \file{agency-blog-saas}; both exist in this portfolio
and both are private.

The page's own prose is looser. In the article reader and in the README, davonex.com is
called ``an actual configured client of the real production pipeline'', without saying
which pipeline. A visitor reading that sentence inside a demo titled ``Agency Blog SaaS''
will very reasonably conclude that these fifteen posts were produced by Agency Blog SaaS.
\textbf{The one file that names a system says otherwise.}

This is worth resolving before an interview, because it is the kind of thing that gets
asked. Two readings are possible and the code cannot settle it: either the two systems
share the client and the wording is merely imprecise, or the library is showcasing another
project's output inside this project's demo. Whichever is true, the honest framing that the
rest of the project works so hard for is undercut by the one ambiguous noun.
\end{honest}

\subsection{The labelling machinery, concretely}

Four separate mechanisms in the code exist only to keep the tiers legible:

\begin{enumerate}
\item \textbf{A permanent banner.} \file{index.html} opens with a fixed red strip:
      \code{Demo mode: fabricated product tour. No live API calls happen here; nothing
      shown is a real client's data.} It is the first element in the body and it never
      goes away.
\item \textbf{Self-labelling data.} The invented client's name is literally stored as
      \code{"Aurora Fitness Co. (invented, illustrative client)"}. The label is part of the
      value, so it cannot be lost by a template that forgets to render a flag.
\item \textbf{The ribbon and countdown.} Any generated card wears a ribbon reading
      \code{Demo}, \code{duplicate}, and \code{removes in 14:59} (joined by middle dots),
      ticking down live.
\item \textbf{The provenance note.} Every opened article ends with a paragraph stating the
      content is real, published, and stored locally, followed by a link to the original
      \file{davonex.com} URL. Verified in a real browser: the link resolves to
      \code{https://davonex.com/blog/<slug>/} for the post that is open.
\end{enumerate}

\begin{honest}{The banner claim ``nothing shown is a real client's data'' is false}
The banner says: \emph{nothing shown is a real client's data}. But the Blog Library is
fifteen real posts from a real client, and the demo's own README, CREDITS and article
provenance notes all say so proudly.

The intended meaning is clearly ``no private client data'', and everything shown is
genuinely public. But as written, the page's most prominent honesty statement contradicts
the page's most interesting honest feature. The footer has the same problem: ``It contains
no real client data''. A careful reader who notices this is being handed a reason to
distrust every other label on the page, which is the exact opposite of what the banner is
for.

This is cheap to fix and worth fixing: ``no private client data; the Blog Library is a real
client's already-public content, labelled as such'' would be both accurate and a stronger
claim.
\end{honest}

\section{Architecture}

\subsection{How the page assembles itself}

There is no framework, so the load sequence is worth tracing exactly once, because
everything else follows from it.

\begin{figure}[H]
\centering
\resizebox{\textwidth}{!}{
\begin{tikzpicture}

\node[box, text width=2.6cm, minimum height=1.5cm] (html) at (0,0)
  {\file{index.html}\\\scriptsize five empty\\ panels};
\node[box, text width=3.5cm, minimum height=1.5cm] (data) at (6.4,0)
  {\file{demo-data.js}\\[2pt]\scriptsize \code{REAL\_DAVONEX\_POSTS}\\
   \code{GENERATOR\_PRESETS}\\ \code{DEMO}};
\node[box, text width=2.7cm, minimum height=1.5cm] (app) at (12.6,0)
  {\file{app.js}\\\scriptsize functions\\ + listeners};
\node[store, text width=1.5cm] (dom) at (17.8,0) {live\\ DOM};

\draw[flow] (html) -- node[lbl, above=1pt] {\code{<script>}\\ tag 1} (data);
\draw[flow] (data) -- node[lbl, above=1pt] {\code{<script>}\\ tag 2\\ (order matters)} (app);
\draw[flow] (app) -- node[lbl, above=1pt] {\code{DOMContentLoaded}\\ fires the\\ 4 renderers} (dom);

\node[extbox, text width=4.4cm] (globals) at (6.4,-3.6)
  {\scriptsize Both files run at top level, so\\ their \code{const} declarations land in\\
   the same global scope. \file{app.js}\\ reads \file{demo-data.js}'s constants\\ by name,
   with no import and no\\ module system.};
\draw[dflow] (globals.north) -- (data.south);

\node[extbox, text width=4.4cm] (renders) at (13.6,-3.6)
  {\scriptsize \code{renderCreditsBanner()}\\ \code{renderPackages()}\\
   \code{renderLibrary()}\\ \code{renderClientUpdates()}\\ each writes \code{.innerHTML} once.};
\draw[dflow] (renders.north) -- (app.south);

\node[box, draw=warnred, fill=warnred!8, text width=4.4cm] (net) at (0,-3.6)
  {\textbf{No network layer exists.}\\\scriptsize There is no \code{fetch}, no
   \code{XMLHttpRequest},\\ no WebSocket, and no external\\ URL loaded anywhere in the page.};
\end{tikzpicture}}
\caption{The complete architecture. Two script tags in a fixed order, three global
constants, four render functions, one DOM.}
\end{figure}

\begin{jargon}{The DOM, and \code{innerHTML}}
The DOM (Document Object Model) is the browser's live, in-memory tree of the page. HTML is
the text you wrote; the DOM is the object the browser built from it and actually draws.
JavaScript changes the page by changing the DOM.

\code{element.innerHTML = "<p>hi</p>"} is the blunt instrument for that: it throws away
everything inside \code{element} and rebuilds it by parsing a string of HTML. It is simple
and fast to write. Its two costs matter later in this document: it destroys and recreates
the elements every time (section~\ref{sec:leak}), and it will execute whatever markup the
string happens to contain (section~\ref{sec:escaping}).
\end{jargon}

\begin{jargon}{Global scope, and why script order is load-bearing}
A browser runs \code{<script>} tags in the order they appear. A \code{const} declared at
the top level of a plain (non-module) script is visible to every script that runs after it.

So \file{app.js}'s very first executable line, \code{let creditsBalance = DEMO.credits;},
works only because \file{demo-data.js} appeared first in \file{index.html} and left
\code{DEMO} lying in the global scope. Swap the two \code{<script>} lines and the page dies
instantly with \code{DEMO is not defined}. This is a real coupling, invisible in either
file, recorded only by the order of two lines in the HTML. It is completely standard for a
page this size, and it is exactly what module systems were invented to make explicit.
\end{jargon}

\subsection{The data model}

\file{demo-data.js} defines three globals and nothing else.

\begin{lstlisting}[caption={The shape of one entry in REAL\_DAVONEX\_POSTS. The body\_html field is elided here; in the real file it is a single string of up to ~20 KB of article markup.}]
const REAL_DAVONEX_POSTS = [
  {
    slug:      "track-ai-search-sales-e-commerce-attribution-guide-2026",
    title:     "Track AI Search Sales: E-Commerce Attribution Guide (2026)",
    excerpt:   "Learn how to track AI search sales for e-commerce. ...",
    by:        "By Davonex",
    date:      "Updated 6 Jul 2026",
    read_time: "5 min read",
    image:     "assets/posts/track-ai-search-sales-...-2026.webp",
    body_html: "<p>To track real e-commerce sales from AI search, ...</p>...",
    references: [
      { title: "A Practical Framework for ...", url: "https://...", date: "2025-12-04" },
      // ... 10 more
    ]
  },
  // ... 14 more posts
];
\end{lstlisting}

\begin{jargon}{Slug}
A short, lowercase, hyphenated, URL-safe version of a title. \code{Track AI Search Sales:
E-Commerce Attribution Guide (2026)} becomes
\code{track-ai-search-sales-e-commerce-attribution-guide-2026}. Here the slug does double
duty: it is the identity of a post inside the demo (every lookup is
\code{REAL\_DAVONEX\_POSTS.find(p => p.slug === slug)}) \emph{and} it is the real path
component on the real site, which is what lets the provenance link be built by string
concatenation rather than stored.
\end{jargon}

The second constant is derived rather than written, and that derivation is the single most
deliberate design decision in the codebase:

\begin{lstlisting}[caption={GENERATOR\_PRESETS is computed from the posts, not authored}]
const GENERATOR_PRESETS = REAL_DAVONEX_POSTS.map(p => ({
  slug: p.slug,
  topic: p.title,
  description: p.excerpt,
  keywords: p.title.toLowerCase().replace(/[():]/g, "")
              .split(" ").filter(w => w.length > 3).slice(0, 6).join(", ")
}));
\end{lstlisting}

\begin{keyidea}{Why derive the presets instead of writing them}
A demo generator has an honesty trap in it. If ``Autofill'' writes one thing into the form
and ``Generate'' produces something unrelated, the demo is quietly lying about the most
basic property a generator has: that the output follows from the input.

By computing each preset \emph{from} the post it will produce, the project makes the link
structural rather than clerical. Preset $i$ carries post $i$'s slug. Autofill records that
slug. Generate looks it up. The same input always yields the same output because they are
the same object, and no future edit can drift them apart, because there is nothing to keep
in sync.

The README states that this mirrors the real pipeline's actual deterministic-per-brief
behaviour. That claim is about the private system and \textbf{cannot be checked from this
repository}; treat it as the owner's assertion, not as something this demo demonstrates.
\end{keyidea}

The third constant holds the fabricated tour data:

\begin{lstlisting}[caption={DEMO, in full}]
const DEMO = {
  credits: 1200,
  packages: [
    { name: "Basic",      price: "$9.99/mo",  credits: "500 credits",   note: "..." },
    { name: "Pro",        price: "$24.99/mo", credits: "1,500 credits", note: "...",
      featured: true },
    { name: "Enterprise", price: "$79.99/mo", credits: "5,000 credits", note: "..." }
  ],
  client: {
    name: "Aurora Fitness Co. (invented, illustrative client)",
    updates: [ { date: "2026-06-01", text: "Published 4 blog posts this month, ..." },
               { date: "2026-05-15", text: "Organic traffic up 12% month over month ..." },
               { date: "2026-05-01", text: "Completed keyword research refresh ..." } ]
  },
  facts: {
    tests:    "132/132 tests passing",
    coverage: "56% coverage",
    stack:    "Django 4.2, PostgreSQL 18, Docker, Gunicorn, nginx"
  }
};
\end{lstlisting}

\begin{honest}{\code{DEMO.facts} is dead code, and its contents are duplicated by hand}
Verified by searching the whole repository: \code{DEMO.facts} is \textbf{never read}. Not by
\file{app.js}, not by \file{index.html}. It is defined and abandoned.

Worse than merely dead: the same three facts are hard-coded as literal text in the footer of
\file{index.html}, which prints ``Built with Django 4.2, PostgreSQL 18, Docker, Gunicorn,
and nginx'' and ``132/132 tests passing, 56\% coverage, per the real project's README''. So
the project has two copies of the same figures, one of which is wired to nothing. Change the
data object and the page does not move; change the page and the data object silently lies.

The fix is one line in each direction and it should be done. As it stands this is a small
but genuine defect in a project whose entire premise is that its claims are traceable.
\end{honest}

\subsection{The claimed figures, and what could be checked}

The footer's ``132/132 tests passing, 56\% coverage'' is careful in one respect: it says
\emph{per the real project's README}, attributing the number rather than asserting it
directly.

\begin{honest}{Verified by cross-reference, not by this repository}
These numbers are not checkable from inside this demo, which has no tests of its own and no
access to the real system. Cross-referencing the private \file{agency-blog-saas} project's
own README confirms that it does state 132 tests and 56\% line coverage against a 50\% gate,
and that it records them as locally verified.

So the demo is quoting its source accurately. What the demo cannot do, and what this
document will not pretend it does, is \emph{verify} them. A reader of the demo has the
owner's word, one step removed. That is the honest status of that footer line, and it is
worth saying out loud in an interview before someone else asks.
\end{honest}

\section{\file{app.js}: every behaviour on the page}

216 lines, no dependencies. This section covers all of it.

\subsection{The state}

The entire application state is five module-level variables.

\begin{lstlisting}[caption={The whole of the demo's memory}]
let creditsBalance = DEMO.credits;   // 1200, decremented by 100 per generate
let libraryDemo = null;              // { slug, addedAt, expiresAt, tickHandle,
                                     //   timeoutHandle } or null
let currentPresetSlug = null;        // set by autofill, read by generate
let readerReturnTab = "library";     // where the Back button goes
const LIBRARY_LIFETIME_MS = 15 * 60 * 1000;   // 15 minutes
\end{lstlisting}

Note what is \emph{not} here: no array of generated posts. \code{libraryDemo} is a single
slot, not a list. Generating twice replaces rather than appends. This is a deliberate
constraint that keeps the library from filling up with duplicates, and it is enforced in
\code{addToLibrary}.

\subsection{The tab machine}

\begin{figure}[H]
\centering
\resizebox{0.96\textwidth}{!}{
\begin{tikzpicture}[node distance=13mm]
\node[box] (gen) {panel-generator\\\scriptsize \emph{active at load}};
\node[box, right=of gen] (lib) {panel-library};
\node[box, right=of lib] (sub) {panel-subscriptions};
\node[box, right=of sub] (cli) {panel-clients};
\node[box, draw=warnred, fill=warnred!8, below=16mm of lib, text width=3.4cm] (read)
  {panel-reader\\\scriptsize \textbf{no tab button exists}};

\draw[flow, <->] (gen) -- (lib);
\draw[flow, <->] (lib) -- (sub);
\draw[flow, <->] (sub) -- (cli);
\draw[flow, <->, bend left=20] (gen) to (sub);
\draw[flow, <->, bend left=34] (gen) to (cli);

\draw[flow] (lib.south) -- node[lbl, right] {click any\\ post card} (read.north);
\draw[flow, bend right=20] (read.west) to node[lbl, below] {\code{reader-back} returns to\\
  \code{readerReturnTab}} (gen.south west);

\node[extbox, right=18mm of read, text width=7.4cm]
  {\scriptsize \code{switchTab(id)} toggles the class \code{is-active} on every \code{.tab}\\
   and every \code{.panel}. Because the reader has no button, entering\\ it leaves
   \emph{all four} tabs looking unselected: the only way out is\\ the Back button inside the
   panel.};
\end{tikzpicture}}
\caption{Five panels, four buttons. The reader is a panel that the tab bar does not know
about.}
\end{figure}

\begin{lstlisting}[caption={The entire router}]
function switchTab(tabId) {
  document.querySelectorAll(".tab")
    .forEach(t => t.classList.toggle("is-active", t.dataset.tab === tabId));
  document.querySelectorAll(".panel")
    .forEach(p => p.classList.toggle("is-active", p.id === `panel-${tabId}`));
}
\end{lstlisting}

Two lines, and it is the right two lines for this problem. \code{classList.toggle(name,
condition)} adds the class when the condition is true and removes it when false, so a single
pass over every tab and every panel puts the whole UI into a consistent state with no
bookkeeping.

\begin{honest}{No URL, no history, no deep link}
The tab is not in the address bar. Refreshing always returns to the generator; the browser
Back button leaves the page entirely rather than leaving the reader; and no article can be
linked to. For a four-tab tour this is a defensible simplification, and adding
\code{history.pushState} would be perhaps fifteen lines. But it does mean an interviewer
cannot be sent a link to a specific screen, which for a demo whose purpose is to be shown
to people is a slightly awkward limitation.
\end{honest}

\subsection{Autofill}

\begin{lstlisting}[caption={autofillRandom}]
function autofillRandom() {
  const preset = GENERATOR_PRESETS[Math.floor(Math.random() * GENERATOR_PRESETS.length)];
  currentPresetSlug = preset.slug;
  document.getElementById("topic-input").value = preset.topic;
  document.getElementById("desc-input").value = preset.description;
  document.getElementById("keywords-input").value = preset.keywords;
  document.getElementById("autofill-note").innerHTML =
    `Filled with a real davonex.com post: <strong>${escapeHtml(preset.topic)}</strong>. ...`;
}
\end{lstlisting}

\code{Math.random()} returns a number in $[0, 1)$; multiplying by the length and flooring
gives a uniformly chosen index. The chosen slug is stashed in \code{currentPresetSlug},
which is the only channel between this function and the generator.

\begin{honest}{The derived keywords are visibly poor, and they are on screen}
The keyword derivation is a title, lowercased, stripped of only \code{(} \code{)} and
\code{:}, split on spaces, filtered to words longer than three characters, and truncated to
six. Running it over all fifteen real titles produces, among others:

\begin{lstlisting}
google, best, e-commerce, 2026            <- from "Google Ads vs AEO: Best E-Commerce ROI in 2026"
what, answer, engine, optimization, e-commerce, 2026?
chatgpt, recommend, your, e-commerce, store, 2026
paid, e-commerce, lowering, 2026          <- from "SEO vs Paid Ads for E-Commerce: Lowering CAC in 2026"
\end{lstlisting}

Three concrete failures, all visible to any visitor who clicks Autofill:
\begin{itemize}
\item \textbf{The length filter deletes the subject.} ``Ads'', ``SEO'', ``AEO'', ``CAC'',
      ``ROI'' and ``vs'' are all four characters or fewer, so the Google Ads post's keywords
      contain neither ``ads'' nor ``AEO'', and the CAC post's contain neither ``SEO'' nor
      ``CAC''. The keywords omit precisely the terms the article is about.
\item \textbf{Stopwords survive.} ``what'', ``your'', ``best'' and ``guide'' are longer than
      three characters, so they stay. No real keyword field would contain ``your''.
\item \textbf{The punctuation strip is incomplete.} It removes \code{()} and \code{:} but
      not \code{?}, so one preset ships the literal keyword \code{2026?}. Confirmed on
      screen in a real browser.
\end{itemize}

This matters more than a normal cosmetic bug, because this demo's pitch is that its
generator screen is a faithful picture of a real product. A keyword field reading
\code{what, answer, engine, optimization, e-commerce, 2026?} is not a picture of a real
product's keyword field; it is a picture of a naive string split. Since the presets are
derived rather than authored, the cheapest honest fix is a small stopword list and a
hand-written keyword line per post, which is fifteen short strings.
\end{honest}

\subsection{The scripted generation}

This is the centrepiece interaction, and it is four timers.

\begin{figure}[H]
\centering
\resizebox{\textwidth}{!}{
\begin{tikzpicture}[x=1.4cm]

\draw[->, thick, draw=inkblue] (0,0) -- (13.4,0) node[right, font=\small] {time};
\foreach \x/\t in {0/0, 2.6/700ms, 7.0/1900ms, 10.7/2900ms}
  { \draw[draw=inkblue] (\x,-0.12) -- (\x,0.12); \node[font=\scriptsize, below=2pt] at (\x,-0.12) {\t}; }

\node[box, above=6mm of {(0,0)}, anchor=south, text width=2.5cm] (s0)
  {\scriptsize \textbf{submit}\\ \code{preventDefault()}\\ button disabled\\ slug captured};
\node[box, above=6mm of {(2.6,0)}, anchor=south, text width=3.2cm] (s1)
  {\scriptsize \textbf{timer 1}\\ \code{creditsBalance -= 100}\\ banner redrawn\\
   ``Generating post...''};
\node[box, above=6mm of {(7.0,0)}, anchor=south, text width=2.6cm] (s2)
  {\scriptsize \textbf{timer 2}\\ ``Formatting and\\ saving to your\\ library...''};
\node[box, above=6mm of {(10.7,0)}, anchor=south, text width=3.2cm] (s3)
  {\scriptsize \textbf{timer 3}\\ \code{addToLibrary(slug)}\\ ``Done: <title>''\\
   button re-enabled};

\node[extbox, below=9mm of {(0,0)}, anchor=north, text width=3.1cm]
  {\scriptsize ``Checking credit\\ balance...'' shown\\ immediately, with\\ a CSS spinner};
\node[extbox, below=9mm of {(6.8,0)}, anchor=north, text width=7.6cm]
  {\scriptsize All three timers are scheduled \emph{at submit}, in parallel, with absolute
   delays. They are not chained. Nothing can fail, so there is no error path, no retry and
   no cancel: the sequence always completes in exactly 2.9 seconds.};
\end{tikzpicture}}
\caption{\code{runGenerator}: a fixed 2.9-second choreography of a process that does not
happen.}
\end{figure}

\begin{jargon}{\code{setTimeout} and \code{preventDefault}}
\code{setTimeout(fn, ms)} asks the browser to call \code{fn} once, after at least \code{ms}
milliseconds, and returns immediately. It does not pause anything.

\code{e.preventDefault()} on a form's \code{submit} event stops the browser's built-in
behaviour, which would be to send the form to a server and reload the page. Since there is
no server, this one line is what makes the whole demo possible: the page stays put and
JavaScript handles the event instead.
\end{jargon}

\begin{honest}{Generating without autofilling gives you someone else's post}
\code{const slug = currentPresetSlug || GENERATOR\_PRESETS[0].slug;}

If a visitor ignores the Autofill button, types their own topic, and hits Generate, then
\code{currentPresetSlug} is still \code{null} and the fallback fires. Verified in a real
browser: typing \code{My Own Unrelated Topic} and generating produces the status line
\code{Done: Track AI Search Sales: E-Commerce Attribution Guide (2026) added to your Blog
Library}.

The form is fully editable, none of the fields are \code{required} or \code{readonly}, and
nothing tells the visitor that their input is being discarded. So the most natural first
interaction on the page, typing something into the topic box, silently produces an
unrelated article and asserts it as the result. That is the demo doing exactly the thing
the preset-derivation design was built to prevent, one branch away from all that care.

The honest options are to disable Generate until Autofill has been clicked, or to make the
inputs read-only with a note, or simply to say in the status line that the demo can only
produce its fifteen stored posts. Any of the three costs a few lines. The current fallback
is the one place where this demo says something untrue without labelling it.
\end{honest}

\begin{honest}{Credits go negative}
\code{creditsBalance -= 100} has no floor and no check. The starting balance is 1{,}200, so
the thirteenth generation takes the display to $-100$, and it keeps going. Verified in a
browser by driving the balance down: the banner will happily render \code{-50}.

In the real system this is the exact place where the interesting logic lives (refuse,
refund on failure, race with a concurrent request). Here it is one subtraction. A visitor
who clicks the button thirteen times sees a negative credit balance on a page that is
supposed to be showing them how a metered product behaves. A single \code{if
(creditsBalance < 100)} guard with an ``out of credits'' message would turn a bug into a
demonstration of the feature.
\end{honest}

\subsection{The demo card's life cycle}
\label{sec:leak}

\begin{figure}[H]
\centering
\resizebox{0.98\textwidth}{!}{
\begin{tikzpicture}[node distance=15mm]
\node[box, text width=3.0cm, minimum height=1.5cm] (none) at (0,0)
  {\code{libraryDemo}\\ \code{= null}\\[2pt]\scriptsize 15 real cards};
\node[box, text width=3.4cm, minimum height=1.5cm] (live) at (9.0,0)
  {\code{libraryDemo = \{...\}}\\[2pt]\scriptsize 16 cards, ribbon\\ counting down};
\node[dec, text width=2.0cm] (again) at (15.4,0) {generate\\ again?};

\draw[flow] (none) -- node[lbl, above=2pt] {\code{addToLibrary(slug)}\\ sets
  \code{expiresAt = now + 15min},\\ starts a 1s \code{setInterval},\\ arms a 15min
  \code{setTimeout}} (live);
\draw[flow] (live) -- (again);
\draw[flow, bend right=32] (again.north) to node[lbl, above=1pt] {yes: old handles cleared,\\
  slot overwritten (never appended)} (live.north);

\draw[flow] (again.south) -- (15.4,-2.4) -- (0,-2.4) -- (none.south);
\node[lbl] at (10.8,-1.9) {15 minutes elapse: \code{setTimeout} fires,\\
  \code{libraryDemo = null}, \code{renderLibrary()}};

\node[box, draw=warnred, fill=warnred!8, text width=8.6cm] (bug) at (6.4,-5.0)
  {\textbf{The expiry path never clears the interval.}\\[3pt]
   \scriptsize \code{timeoutHandle: setTimeout(() => \{ libraryDemo = null;
   renderLibrary(); \}, ...)}\\[3pt]
   \scriptsize Nulling \code{libraryDemo} discards the only reference to
   \code{tickHandle}, so \code{clearInterval} can never be called for it again. The 1s
   ticker runs for the life of the page.};
\draw[dflow] (bug.north) -- (6.4,-2.4);
\end{tikzpicture}}
\caption{The single-slot life cycle, and the leak on the expiry edge.}
\end{figure}

\begin{lstlisting}[caption={addToLibrary, in full}]
function addToLibrary(slug) {
  if (libraryDemo) {                        // replacing: clean up the old one
    clearInterval(libraryDemo.tickHandle);
    clearTimeout(libraryDemo.timeoutHandle);
  }
  const addedAt = Date.now();
  const expiresAt = addedAt + LIBRARY_LIFETIME_MS;
  libraryDemo = {
    slug, addedAt, expiresAt,
    tickHandle: setInterval(renderLibrary, 1000),
    timeoutHandle: setTimeout(() => { libraryDemo = null; renderLibrary(); },
                              LIBRARY_LIFETIME_MS)
  };
  renderLibrary();
}
\end{lstlisting}

The replace path is careful and correct: it clears both handles before overwriting the slot.
The expiry path is not.

\begin{honest}{Verified leak: the ticker outlives the card, forever}
Driven in a real headless Chrome. A \code{MutationObserver} was attached to
\code{\#library-list} to count rebuilds, then the expiry callback was executed exactly as
\file{app.js} writes it.

\begin{lstlisting}
DOM rebuilds of the 16-card grid in 3s while the demo card is alive:  4
DOM rebuilds in 3s AFTER the demo card expired:                       4   <-- still ticking
\end{lstlisting}

The interval keeps firing at 1\,Hz forever, and each fire runs \code{renderLibrary()}, which
rebuilds the \emph{entire fifteen-card grid} via \code{innerHTML}: fifteen
\code{<article>} elements, fifteen \code{<img>} tags, all destroyed and re-parsed, once a
second, permanently, for a demo card that no longer exists. The handle is unreachable, so
nothing can ever stop it short of a page reload.

The fix is one line: capture the handle in a local variable, or call
\code{clearInterval(libraryDemo.tickHandle)} inside the timeout callback before nulling the
slot.

Note the second-order cost too. Even \emph{before} expiry, this design rebuilds fifteen
image elements every second purely to update the \code{mm:ss} in one ribbon. Updating a
single text node instead of re-rendering the grid would be both correct and cheaper. This
is the clearest ``I would do it differently'' answer the codebase offers, and it is worth
having ready.
\end{honest}

\subsection{The article reader}

\code{openReader(slug, isDemo)} builds an entire article as one \code{innerHTML} string:
hero image, an optional demo flag, title, byline, the stored \code{body\_html} verbatim, a
references list, and the provenance note with its link back to the original.

One small nicety before it renders:

\begin{lstlisting}
readerReturnTab = document.querySelector(".panel.is-active")?.id.replace("panel-", "")
                  || "library";
\end{lstlisting}

The reader remembers which panel it was opened from so Back returns there, defaulting to the
library. Verified: Back does return to \code{panel-library}. In practice only the library
has clickable cards, so the generality is unused, but it costs one line and it is correct.

\begin{jargon}{Optional chaining, \code{?.}}
\code{a?.b} evaluates to \code{undefined} instead of throwing if \code{a} is \code{null} or
\code{undefined}. Here it guards the case where no panel carries \code{is-active}, which
then falls through to the \code{|| "library"} default. It is a defensive habit; the case
cannot actually arise, because a panel is always active.
\end{jargon}

\begin{honest}{Every inline citation link in every article is dead}
The real articles are full of citation markers: \code{<sup class="citation"><a
href="\#ref-1">[1]</a></sup>}. They are meant to jump to the matching entry in the
references list.

But \code{openReader} renders the references as a plain \code{<ol>} of \code{<li>} elements
with \textbf{no \code{id} attributes}, and searching all fifteen stored bodies confirms that
no \code{id="ref-N"} target exists anywhere in the data either. So every one of those
anchors points at an element that does not exist.

Measured in a real browser on one opened article: \textbf{21 citation links, 21 of them
dead.} Clicking any of them does nothing at all.

This is the price of copying rendered HTML out of a site and re-hosting it in a different
template: the markup came across, the anchor targets it depended on did not. It is a
one-line fix (\code{<li id="ref-\$\{i+1\}">} with an index from \code{map}), and it is
worth doing precisely because these citations are the part of the demo that shows real,
sourced work. A panel member who clicks a citation and gets nothing has just found the
seam.
\end{honest}

\subsection{Escaping, and the one place it is skipped}
\label{sec:escaping}

\begin{lstlisting}[caption={The project's single security primitive}]
function escapeHtml(str) {
  return String(str).replace(/&/g, "&amp;").replace(/</g, "&lt;").replace(/>/g, "&gt;");
}
\end{lstlisting}

\begin{jargon}{Why escaping exists (XSS)}
When you build HTML by pasting a string into a template, the browser cannot tell your markup
from the pasted string's markup. If the string contains \code{<script>}, the browser runs it.
This is cross-site scripting (XSS).

Escaping neutralises the characters that start markup: \code{<} becomes \code{\&lt;}, which
the browser draws as a less-than sign rather than treating as the start of a tag. Escaping
\code{\&} first is essential and correctly ordered here; doing it last would double-escape
the entities the earlier replacements just introduced.
\end{jargon}

The function is applied consistently to every title, excerpt, date, price and update text.
Two deliberate exceptions and one apparent oversight:

\begin{itemize}
\item \code{post.body\_html} is injected raw. This is intentional and unavoidable: it
      \emph{is} HTML, and escaping it would print the tags as text.
\item \code{post.image} is interpolated into a \code{src} unescaped. It is a local path from
      a static file, so nothing can reach it.
\item \code{r.url} is interpolated into an \code{href} unescaped, in the references list.
\end{itemize}

\begin{honest}{The escaping is theatre, and it is worth knowing that}
Every string \code{escapeHtml} touches is a hard-coded literal in \file{demo-data.js}. There
is no user input anywhere on this page that reaches the DOM: the three form fields are read
into \code{currentPresetSlug} logic and the redeem box is only tested for emptiness. Its
value is never rendered.

So \code{escapeHtml} defends against nothing that can happen here. The unescaped \code{href}
is equally harmless for the same reason: an attacker would have to edit the repository to
exploit it, at which point they can do anything anyway.

This is a good habit applied to a codebase that does not need it, which is the right way
round. But it should be described accurately: the demo is safe because it has \emph{no
untrusted input}, not because it escapes carefully. The two are different arguments and only
the first one is doing any work.
\end{honest}

\subsection{Wiring}

\begin{lstlisting}[caption={DOMContentLoaded: the only entry point}]
document.addEventListener("DOMContentLoaded", () => {
  renderCreditsBanner(); renderPackages(); renderLibrary(); renderClientUpdates();

  document.querySelectorAll(".tab")
    .forEach(t => t.addEventListener("click", () => switchTab(t.dataset.tab)));
  document.getElementById("generator-form").addEventListener("submit", runGenerator);
  document.getElementById("autofill-btn").addEventListener("click", autofillRandom);

  document.getElementById("library-list").addEventListener("click", (e) => {
    const card = e.target.closest("[data-open-post]");
    if (!card) return;
    openReader(card.dataset.openPost, card.dataset.isDemo === "1");
  });
  document.getElementById("reader-back").addEventListener("click", () =>
    switchTab(readerReturnTab));
  // ... redeem handler
});
\end{lstlisting}

\begin{keyidea}{Event delegation: the one non-obvious technique in the file}
The library's cards are destroyed and rebuilt every second while a countdown runs. A
listener attached to a card would die with the card.

So no listener is attached to any card. One listener sits on the container,
\code{\#library-list}, which is never replaced. When a click arrives, \code{e.target} is
whatever exact element was hit (possibly the image, or the title text), and
\code{e.target.closest("[data-open-post]")} walks \emph{up} the tree to find the nearest
ancestor carrying that attribute, which is the card. If there is none, the click was on
empty space and is ignored.

This is called event delegation. It works because events bubble up the DOM tree, and it is
the correct answer to the ``my listeners disappear when I re-render'' problem. The card's
identity travels in \code{data-open-post} (the slug) and \code{data-is-demo}, read back via
\code{dataset}, so the handler needs no closure over anything.
\end{keyidea}

\section{The self-contained guarantee}
\label{sec:selfcontained}

The demo makes a strong, testable promise, repeated in the banner, the footer, the README,
and the first comment line of both JavaScript files: \textbf{no network calls}.

\begin{keyidea}{Why this is a real design constraint and not a slogan}
An honesty claim that depends on a remote server is not an honesty claim. If the demo
fetched davonex.com's posts live, then: the page would break when the client's site changed
or went down; the demo would be reporting on a live system it does not control; and ``no
live API calls'' would be false. Every article body, every hero image and the font itself
were therefore fetched once (dated 2026-07-07 in \file{CREDITS.md}) and committed, which is
why the repository is 2.5\,MB for 1{,}200 lines of code.

The README is explicit that this is the constraint it refused to break, and names the cost:
the fifteen-post snapshot will drift as the real client publishes more, and re-syncing was
skipped rather than adding ``the one exception that would break it''. That is a
proportionate trade, honestly stated.
\end{keyidea}

\begin{keyidea}{Verified, not assumed}
The page was loaded in a real headless Chrome from a \code{file://} origin with every
request intercepted and logged.

\textbf{Result: 0 requests to any non-local origin.} No HTTP, no fetch, no beacon, no
analytics, no font CDN, no hotlinked images. The page also produced \textbf{zero JavaScript
errors} across the full interaction: load, autofill, generate, open an article, go back,
generate again, and expire the demo card.

The claim is true. This is the one headline claim the demo makes that this document can
confirm from the outside rather than take on trust.
\end{keyidea}

The font is part of that: \file{assets/fonts/space-grotesk.woff2} is committed and declared
with a local \code{@font-face} rather than pulled from Google Fonts, which would have been a
network call and a third-party request on every load.

\section{\file{styles.css}: a borrowed design system}

285 lines, opening with a comment that is itself a piece of the project's honesty
discipline:

\begin{lstlisting}[caption={The header comment of styles.css, quoted}]
  Palette and component patterns below match the real, live public design
  system read directly off davonex.com's own CSS (:root tokens, .card/.btn/
  .eyebrow components), per the owner's direction to style this demo after
  davonex.com rather than digitalise.agency's own indigo palette. Colors,
  radii, and shadows are copied 1:1; nothing here is invented. This is public
  visual design language (what every visitor's browser already renders), not
  any site's private server code, auth, or billing logic, which stays out of
  this repo entirely.
\end{lstlisting}

The tokens themselves are CSS custom properties:

\begin{lstlisting}[caption={A representative slice of the :root block}]
:root {
  --bg: #FAFBFE;          --paper: #FFFFFF;
  --blue: #136BEE;        --blue-deep: #0C4497;    --navy: #062757;
  --ink: #08152B;         --ink-soft: #46597A;
  --line: rgba(6, 39, 87, .12);
  --r-lg: 26px;           --r-md: 16px;            --r-sm: 10px;
  --shadow-blue: 0 18px 40px rgba(19, 107, 238, .28);
  --wrap: 960px;
}
\end{lstlisting}

\begin{jargon}{CSS custom properties (variables)}
\code{--blue: \#136BEE} declares a named value on an element; \code{var(--blue)} reads it,
and it cascades to all descendants. Declaring them on \code{:root} (the \code{<html>}
element) makes them global. The point is that the palette lives in one block rather than
being scattered as literal hex codes across 285 lines, so a token can be re-pointed once.
\end{jargon}

The rest of the file is unremarkable modern CSS, and readable: a fixed banner, a sticky
blurred nav, a card with a hover lift and a top accent bar, arrow-animated buttons, a
keyframed spinner, a grid for the library, and an article template mirroring the real site's
post layout. There is no preprocessor and no framework, so what is in the file is what the
browser gets.

\begin{honest}{The provenance argument is stated but not sourced}
The comment says the values were ``read directly off davonex.com's own CSS'' and copied
1:1. That is a specific, checkable claim, and it cannot be checked from this repository:
there is no record of which stylesheet, at what date, and the demo makes no live request
that could confirm it. \file{CREDITS.md} carefully dates the article snapshot (2026-07-07)
but says nothing about the CSS.

The reasoning offered (that a public site's rendered visual language is not private code) is
reasonable, and the comment is admirably specific about the owner's direction. But the
project's own standard elsewhere is to date and source its borrowings, and here it does not.
Adding a line to \file{CREDITS.md} would close the gap.
\end{honest}

\section{Licensing}

The project ships PolyForm Strict License 1.0.0, copyright 2026 Salman Adnan. This is a
\emph{source-available} licence, not an open-source one: it permits reading and evaluating
the code, and forbids use, modification and redistribution.

\begin{keyidea}{Why a strict licence is the right call here}
The repository contains a real client's copyrighted article text and images, reproduced with
permission from the owner's own agency relationship rather than under any open licence
(\file{CREDITS.md} says so explicitly). An MIT licence on this repository would purport to
grant strangers rights to fifteen articles that are not the licensor's to give away that
freely. PolyForm Strict is the licence that matches what the repository actually is: a
thing to be looked at, not a thing to be taken.
\end{keyidea}

\section{Everything that was verified, and how}

The brief for this document requires that claims be run rather than assumed wherever that is
cheap. It was cheap here. The page was driven in a real headless Chrome; the data file was
evaluated in Node.

\begin{table}[H]
\centering
\small
\begin{tabular}{@{}p{0.46\textwidth}p{0.20\textwidth}p{0.26\textwidth}@{}}
\toprule
\textbf{Claim} & \textbf{Result} & \textbf{Method} \\
\midrule
No network calls at runtime & \textbf{True} & 0 non-local requests intercepted \\
No JavaScript errors across the full tour & \textbf{True} & 0 pageerror, 0 console errors \\
15 real posts, all present at load & \textbf{True} & 15 \code{.post-card} elements \\
15 unique slugs, 15 unique titles & \textbf{True} & Set comparison in Node \\
All 15 hero images exist on disk & \textbf{True} & \code{fs.existsSync} per post \\
Generate adds exactly one card & \textbf{True} & 15 to 16, label \code{(16)} \\
Credits go 1{,}200 to 1{,}100 per generate & \textbf{True} & read from the DOM \\
Autofilled post is the generated post & \textbf{True} & titles compared, equal \\
Countdown starts at 15 minutes & \textbf{True} & ribbon read: \code{removes in 14:59} \\
Provenance link points at the open post & \textbf{True} & href read from the DOM \\
Back returns to the previous panel & \textbf{True} & \code{panel-library} \\
Article renders with tables and refs & \textbf{True} & 13 paragraphs, 1 table, 11 refs \\
\midrule
Citation links resolve & \textbf{False} & 21 links, 21 dead \\
Typed topic is respected & \textbf{False} & unrelated post returned \\
Credits have a floor & \textbf{False} & rendered $-50$ \\
Ticker stops when the card expires & \textbf{False} & 4 rebuilds per 3s after expiry \\
\code{DEMO.facts} is used & \textbf{False} & no reference in the repository \\
\midrule
132/132 tests, 56\% coverage & \textbf{Not verifiable here} & matches the private
project's README; the demo has no tests of its own \\
Real pipeline is deterministic per brief & \textbf{Not verifiable here} & claim about the
private system \\
CSS tokens copied 1:1 from davonex.com & \textbf{Not verifiable here} & no dated source
recorded \\
\bottomrule
\end{tabular}
\caption{Twelve claims confirmed, five defects confirmed, three claims that this repository
cannot settle either way.}
\end{table}

\section{The honest summary}

\subsection{What is genuinely good here}

\begin{itemize}
\item \textbf{The self-contained guarantee holds.} It is the demo's loudest claim and it
      survives measurement, at zero non-local requests. Most demos that promise this leak a
      font or an analytics beacon; this one does not.
\item \textbf{Deriving the presets from the posts} is a real design decision with a real
      argument behind it, and it is the sort of thing worth being asked about.
\item \textbf{The three-tier labelling} is more thought than most demos get, and shipping a
      real client's public work next to labelled fakes is a genuinely stronger position than
      inventing a library.
\item \textbf{Event delegation on a re-rendering grid} is the correct technique, applied for
      the correct reason.
\item \textbf{The code is small, clean and legible.} Four files, one job each, 216 lines of
      behaviour, no dependencies, no build. For what this needs to do, that is right.
\end{itemize}

\subsection{What does not hold up}

\begin{enumerate}
\item \textbf{The banner's central claim is false as written} (``nothing shown is a real
      client's data''), and it contradicts the project's own best feature.
\item \textbf{The no-autofill path returns an unrelated post and calls it the result},
      which is the one unlabelled untruth on the page.
\item \textbf{Every citation link in every article is dead}: 21 of 21 on the post measured.
\item \textbf{The interval leaks forever} and rebuilds fifteen images per second for the
      life of the page after a card expires.
\item \textbf{Credits go negative}, in a demo whose subject is metering.
\item \textbf{\code{DEMO.facts} is dead}, and its figures are duplicated by hand in the
      footer.
\item \textbf{The keyword derivation drops the subject of the article} and keeps the
      stopwords, on screen, in the flagship interaction.
\item \textbf{Two advertised capabilities} (per-client editor roles, PDF export) are not in
      the tour at all.
\item \textbf{The davonex provenance names a different private system} in \file{CREDITS.md}
      than the demo's prose implies.
\end{enumerate}

None of these is hard to fix. Items 1, 2, 3 and 5 are each a handful of lines, and they are
the four that a careful visitor is most likely to trip over.

\subsection{What this demo cannot demonstrate, and should not be asked to}

It is worth being blunt, because the gap between the demo and the product is where an
interview will go.

This repository contains no evidence about: whether the real generation works, how good its
output is, how prompts are constructed, how credits are actually metered or refunded, how
subscriptions are billed or provisioned, how per-client roles are enforced, how PDF export
is done, how any of it performs, whether it is secure, whether it scales, or whether the
tests it cites exist. It demonstrates that its author can build a small, clean, honest
static page and think carefully about how to label a fiction. That is a real skill and it is
on display here. It is not the same skill as building the Django system, and this demo is
not evidence of that one.

\begin{keyidea}{The strongest defensible claim}
``This is a puppet, I built the puppet carefully, and I labelled it as a puppet.'' Everything
above supports that sentence. Nothing above supports any sentence about the real product's
quality, and the demo is at its best when it does not try to.
\end{keyidea}

\end{document}
